% --- Archon physics formalization source begin ---
% archon:physics
% archon:covers PhyXMiniProblems/problem_phyx_mini_0983.lean
% archon:source-report reports/phyx_mini/problem_phyx_mini_0983.source.json

\chapter{Physics problem phyx\_mini\_0983}
\label{ch:PhyXMiniProblems_problem_phyx_mini_0983}

\paragraph{Problem source.}
\#\# Physical scenario

A uniform electric field is directed axially in a cylindrical region that includes a rectangular loop of wire with total resistance \textbackslash{}( R \textbackslash{}). This loop has radially oriented width \textbackslash{}( a \textbackslash{}) and axially oriented length \textbackslash{}( b \textbackslash{}), and sits tight against the cylinder axis. The electric field is zero at time \textbackslash{}( t = 0 \textbackslash{}) and then increases in time according to \textbackslash{}( \textbackslash{}vec\{E\} = \textbackslash{}eta t\textasciicircum{}2 \textbackslash{}hat\{k\} \textbackslash{}), where \textbackslash{}( \textbackslash{}eta \textbackslash{}) is a constant with units of V/(m\textbackslash{}cdot s\textasciicircum{}2).

\#\# Question

At time \textbackslash{}( t = 5.00 \textbackslash{}text\{ s\} \textbackslash{}), what is the direction of the induced current as viewed from above?

\#\# Answer choices

- A: A: \textbackslash{}frac\{\{\textbackslash{}mu\_0 \textbackslash{}eta b\textasciicircum{}2\}\}\{\{2R\}\}

- B: B: \textbackslash{}frac\{\{\textbackslash{}mu\_0 \textbackslash{}varepsilon\_0 \textbackslash{}eta b\textasciicircum{}2 a\textasciicircum{}2\}\}\{\{2R\}\}

- C: C: \textbackslash{}frac\{\{\textbackslash{}mu\_0 \textbackslash{}varepsilon\_0 \textbackslash{}eta a\textasciicircum{}2\}\}\{\{2R\}\}

- D: D: \textbackslash{}mu\_0 \textbackslash{}varepsilon\_0 \textbackslash{}eta ba

\#\# Figure caption (auxiliary; use the image as primary evidence)

The image depicts a diagram involving a cylindrical shape with a cross-sectional loop inside. The cylinder has a series of pink arrows pointing upward along the central vertical axis, labeled as \textbackslash{}(\textbackslash{}vec\{E\}\textbackslash{}), representing an electric field.

Inside the cylinder, there's a circular loop with radius \textbackslash{}(r\textbackslash{}). The loop is connected to a resistor labeled \textbackslash{}(R\textbackslash{}) in an external circuit, featuring two perpendicular segments labeled \textbackslash{}(a\textbackslash{}) and \textbackslash{}(b\textbackslash{}).

The cylinder and arrows are drawn in pink, and the loop and circuit are in black. The dashed line in the center of the cylinder denotes the axis of symmetry.

\#\# Dataset metadata

Electromagnetic Induction; Physical Model Grounding Reasoning, Spatial Relation Reasoning

\paragraph{Recorded answer/context.}
B: B: \textbackslash{}frac\{\{\textbackslash{}mu\_0 \textbackslash{}varepsilon\_0 \textbackslash{}eta b\textasciicircum{}2 a\textasciicircum{}2\}\}\{\{2R\}\}

\paragraph{Figure/image path.}
/root/proposal\_for\_physic/science-mango/phyx\_mini\_run/phyx\_data/test\_image/983.png

\paragraph{Formalization target.}
create a compiling Lean file with sorry bodies at `PhyXMiniProblems/problem\_phyx\_mini\_0983.lean`. The Lean declarations must preserve the physical quantities, dimensions or dimensional roles, figure labels, governing-law hypotheses, and final relation expressed by this problem.
Use Mathlib/Physlib names found through LeanExplore where available. If a physics API is missing, introduce faithful local abstractions rather than scalar placeholder aliases.

\begin{theorem}[Physics formalization target]
\label{thm:physics:phyx_mini_0983:target}
\lean{PhyXMiniProblems.ProblemPhyXMini0983.problem_phyx_mini_0983}
\uses{def:physics:phyx-mini-0983:phyxminiproblems-problemphyxmini0983-lengthinmeters, def:physics:phyx-mini-0983:phyxminiproblems-problemphyxmini0983-coefficientinvoltspermetersecondsquared, def:physics:phyx-mini-0983:phyxminiproblems-problemphyxmini0983-electriccurrentmagnitudeinamperes, def:physics:phyx-mini-0983:phyxminiproblems-problemphyxmini0983-resistanceinohms, def:physics:phyx-mini-0983:phyxminiproblems-problemphyxmini0983-permittivityinfaradspermeter, def:physics:phyx-mini-0983:phyxminiproblems-problemphyxmini0983-permeabilityinhenriespermeter, def:physics:phyx-mini-0983:phyxminiproblems-problemphyxmini0983-displacementcurrentinductionsetup, def:physics:phyx-mini-0983:phyxminiproblems-problemphyxmini0983-matchesproblemandprimaryfigure, def:physics:phyx-mini-0983:phyxminiproblems-problemphyxmini0983-satisfiesradialaxialcircuitgeometry, def:physics:phyx-mini-0983:phyxminiproblems-problemphyxmini0983-hasphysicalinductionparameters, def:physics:phyx-mini-0983:phyxminiproblems-problemphyxmini0983-satisfiesquadraticaxialelectricfieldmodel, def:physics:phyx-mini-0983:phyxminiproblems-problemphyxmini0983-satisfiescylindricalamperemaxwelllaw, def:physics:phyx-mini-0983:phyxminiproblems-problemphyxmini0983-satisfiesfluxfaradayohmanddirectionlaws}
The assigned autoformalize agent should translate this physics problem into Lean declarations in the covered file, with theorem and lemma proof bodies written as `by sorry`.
\end{theorem}
\begin{proof}
This is an autoformalization task, not a proof task. Produce statements that can later be proved by the physics prover without weakening the physical model.
\end{proof}

% --- Archon declaration topology begin ---
\section{Declaration topology}
\begin{definition}[electric Field Dimension]
  \label{def:physics:phyx-mini-0983:phyxminiproblems-problemphyxmini0983-electricfielddimension}
  \lean{PhyXMiniProblems.ProblemPhyXMini0983.electricFieldDimension}
  Electric-field dimension M L T\textasciicircum{}-2 C\textasciicircum{}-1 (volt per metre)
\end{definition}
\begin{proof}
  This is the declared model definition of electric Field Dimension.
\end{proof}
\begin{definition}[electric Field Quadratic Coefficient Dimension]
  \label{def:physics:phyx-mini-0983:phyxminiproblems-problemphyxmini0983-electricfieldquadraticcoefficientdimension}
  \lean{PhyXMiniProblems.ProblemPhyXMini0983.electricFieldQuadraticCoefficientDimension}
  \uses{def:physics:phyx-mini-0983:phyxminiproblems-problemphyxmini0983-electricfielddimension}
  Dimension of eta in E(t) = eta * t\textasciicircum{}2, namely electric field per time squared
\end{definition}
\begin{proof}
  This is the declared model definition of electric Field Quadratic Coefficient Dimension.
\end{proof}
\begin{definition}[electric Field Rate Dimension]
  \label{def:physics:phyx-mini-0983:phyxminiproblems-problemphyxmini0983-electricfieldratedimension}
  \lean{PhyXMiniProblems.ProblemPhyXMini0983.electricFieldRateDimension}
  \uses{def:physics:phyx-mini-0983:phyxminiproblems-problemphyxmini0983-electricfielddimension}
  Dimension of the time derivative of an electric-field component
\end{definition}
\begin{proof}
  This is the declared model definition of electric Field Rate Dimension.
\end{proof}
\begin{definition}[magnetic Flux Density Dimension]
  \label{def:physics:phyx-mini-0983:phyxminiproblems-problemphyxmini0983-magneticfluxdensitydimension}
  \lean{PhyXMiniProblems.ProblemPhyXMini0983.magneticFluxDensityDimension}
  Magnetic-flux-density dimension M T\textasciicircum{}-1 C\textasciicircum{}-1 (tesla)
\end{definition}
\begin{proof}
  This is the declared model definition of magnetic Flux Density Dimension.
\end{proof}
\begin{definition}[magnetic Flux Dimension]
  \label{def:physics:phyx-mini-0983:phyxminiproblems-problemphyxmini0983-magneticfluxdimension}
  \lean{PhyXMiniProblems.ProblemPhyXMini0983.magneticFluxDimension}
  \uses{def:physics:phyx-mini-0983:phyxminiproblems-problemphyxmini0983-magneticfluxdensitydimension}
  Magnetic-flux dimension M L\textasciicircum{}2 T\textasciicircum{}-1 C\textasciicircum{}-1 (weber)
\end{definition}
\begin{proof}
  This is the declared model definition of magnetic Flux Dimension.
\end{proof}
\begin{definition}[electromotive Force Dimension]
  \label{def:physics:phyx-mini-0983:phyxminiproblems-problemphyxmini0983-electromotiveforcedimension}
  \lean{PhyXMiniProblems.ProblemPhyXMini0983.electromotiveForceDimension}
  \uses{def:physics:phyx-mini-0983:phyxminiproblems-problemphyxmini0983-magneticfluxdimension}
  Magnetic-flux-rate and electromotive-force dimension (volt)
\end{definition}
\begin{proof}
  This is the declared model definition of electromotive Force Dimension.
\end{proof}
\begin{definition}[electric Current Dimension]
  \label{def:physics:phyx-mini-0983:phyxminiproblems-problemphyxmini0983-electriccurrentdimension}
  \lean{PhyXMiniProblems.ProblemPhyXMini0983.electricCurrentDimension}
  Electric-current dimension C T\textasciicircum{}-1 (ampere)
\end{definition}
\begin{proof}
  This is the declared model definition of electric Current Dimension.
\end{proof}
\begin{definition}[electrical Resistance Dimension]
  \label{def:physics:phyx-mini-0983:phyxminiproblems-problemphyxmini0983-electricalresistancedimension}
  \lean{PhyXMiniProblems.ProblemPhyXMini0983.electricalResistanceDimension}
  Electrical-resistance dimension M L\textasciicircum{}2 T\textasciicircum{}-1 C\textasciicircum{}-2 (ohm)
\end{definition}
\begin{proof}
  This is the declared model definition of electrical Resistance Dimension.
\end{proof}
\begin{definition}[vacuum Permittivity Dimension]
  \label{def:physics:phyx-mini-0983:phyxminiproblems-problemphyxmini0983-vacuumpermittivitydimension}
  \lean{PhyXMiniProblems.ProblemPhyXMini0983.vacuumPermittivityDimension}
  Vacuum-permittivity dimension C\textasciicircum{}2 T\textasciicircum{}2 M\textasciicircum{}-1 L\textasciicircum{}-3
\end{definition}
\begin{proof}
  This is the declared model definition of vacuum Permittivity Dimension.
\end{proof}
\begin{definition}[vacuum Permeability Dimension]
  \label{def:physics:phyx-mini-0983:phyxminiproblems-problemphyxmini0983-vacuumpermeabilitydimension}
  \lean{PhyXMiniProblems.ProblemPhyXMini0983.vacuumPermeabilityDimension}
  Vacuum-permeability dimension M L C\textasciicircum{}-2
\end{definition}
\begin{proof}
  This is the declared model definition of vacuum Permeability Dimension.
\end{proof}
\begin{definition}[Time Quantity]
  \label{def:physics:phyx-mini-0983:phyxminiproblems-problemphyxmini0983-timequantity}
  \lean{PhyXMiniProblems.ProblemPhyXMini0983.TimeQuantity}
  A nonnegative, unit-independent physical time
\end{definition}
\begin{proof}
  This is the declared model definition of Time Quantity.
\end{proof}
\begin{definition}[Length Quantity]
  \label{def:physics:phyx-mini-0983:phyxminiproblems-problemphyxmini0983-lengthquantity}
  \lean{PhyXMiniProblems.ProblemPhyXMini0983.LengthQuantity}
  A nonnegative, unit-independent physical length
\end{definition}
\begin{proof}
  This is the declared model definition of Length Quantity.
\end{proof}
\begin{definition}[Electric Field Quadratic Coefficient Magnitude]
  \label{def:physics:phyx-mini-0983:phyxminiproblems-problemphyxmini0983-electricfieldquadraticcoefficientmagnitude}
  \lean{PhyXMiniProblems.ProblemPhyXMini0983.ElectricFieldQuadraticCoefficientMagnitude}
  \uses{def:physics:phyx-mini-0983:phyxminiproblems-problemphyxmini0983-electricfieldquadraticcoefficientdimension}
  The nonnegative physical coefficient eta in E(t) = eta * t\textasciicircum{}2
\end{definition}
\begin{proof}
  This is the declared model definition of Electric Field Quadratic Coefficient Magnitude.
\end{proof}
\begin{definition}[Signed Electric Field Component]
  \label{def:physics:phyx-mini-0983:phyxminiproblems-problemphyxmini0983-signedelectricfieldcomponent}
  \lean{PhyXMiniProblems.ProblemPhyXMini0983.SignedElectricFieldComponent}
  \uses{def:physics:phyx-mini-0983:phyxminiproblems-problemphyxmini0983-electricfielddimension}
  A signed axial electric-field component
\end{definition}
\begin{proof}
  This is the declared model definition of Signed Electric Field Component.
\end{proof}
\begin{definition}[Signed Electric Field Rate]
  \label{def:physics:phyx-mini-0983:phyxminiproblems-problemphyxmini0983-signedelectricfieldrate}
  \lean{PhyXMiniProblems.ProblemPhyXMini0983.SignedElectricFieldRate}
  \uses{def:physics:phyx-mini-0983:phyxminiproblems-problemphyxmini0983-electricfieldratedimension}
  A signed time derivative of an axial electric-field component
\end{definition}
\begin{proof}
  This is the declared model definition of Signed Electric Field Rate.
\end{proof}
\begin{definition}[Signed Magnetic Flux Density Component]
  \label{def:physics:phyx-mini-0983:phyxminiproblems-problemphyxmini0983-signedmagneticfluxdensitycomponent}
  \lean{PhyXMiniProblems.ProblemPhyXMini0983.SignedMagneticFluxDensityComponent}
  \uses{def:physics:phyx-mini-0983:phyxminiproblems-problemphyxmini0983-magneticfluxdensitydimension}
  A signed azimuthal magnetic-flux-density component
\end{definition}
\begin{proof}
  This is the declared model definition of Signed Magnetic Flux Density Component.
\end{proof}
\begin{definition}[Signed Magnetic Flux]
  \label{def:physics:phyx-mini-0983:phyxminiproblems-problemphyxmini0983-signedmagneticflux}
  \lean{PhyXMiniProblems.ProblemPhyXMini0983.SignedMagneticFlux}
  \uses{def:physics:phyx-mini-0983:phyxminiproblems-problemphyxmini0983-magneticfluxdimension}
  Signed magnetic flux through the oriented rectangular circuit
\end{definition}
\begin{proof}
  This is the declared model definition of Signed Magnetic Flux.
\end{proof}
\begin{definition}[Signed Magnetic Flux Rate]
  \label{def:physics:phyx-mini-0983:phyxminiproblems-problemphyxmini0983-signedmagneticfluxrate}
  \lean{PhyXMiniProblems.ProblemPhyXMini0983.SignedMagneticFluxRate}
  \uses{def:physics:phyx-mini-0983:phyxminiproblems-problemphyxmini0983-electromotiveforcedimension}
  Signed magnetic-flux rate at the observation time
\end{definition}
\begin{proof}
  This is the declared model definition of Signed Magnetic Flux Rate.
\end{proof}
\begin{definition}[Electromotive Force]
  \label{def:physics:phyx-mini-0983:phyxminiproblems-problemphyxmini0983-electromotiveforce}
  \lean{PhyXMiniProblems.ProblemPhyXMini0983.ElectromotiveForce}
  \uses{def:physics:phyx-mini-0983:phyxminiproblems-problemphyxmini0983-electromotiveforcedimension}
  Signed electromotive force relative to the chosen circuit traversal
\end{definition}
\begin{proof}
  This is the declared model definition of Electromotive Force.
\end{proof}
\begin{definition}[Electric Current]
  \label{def:physics:phyx-mini-0983:phyxminiproblems-problemphyxmini0983-electriccurrent}
  \lean{PhyXMiniProblems.ProblemPhyXMini0983.ElectricCurrent}
  \uses{def:physics:phyx-mini-0983:phyxminiproblems-problemphyxmini0983-electriccurrentdimension}
  Signed electric current relative to the chosen circuit traversal
\end{definition}
\begin{proof}
  This is the declared model definition of Electric Current.
\end{proof}
\begin{definition}[Electrical Resistance Magnitude]
  \label{def:physics:phyx-mini-0983:phyxminiproblems-problemphyxmini0983-electricalresistancemagnitude}
  \lean{PhyXMiniProblems.ProblemPhyXMini0983.ElectricalResistanceMagnitude}
  \uses{def:physics:phyx-mini-0983:phyxminiproblems-problemphyxmini0983-electricalresistancedimension}
  A nonnegative total circuit resistance
\end{definition}
\begin{proof}
  This is the declared model definition of Electrical Resistance Magnitude.
\end{proof}
\begin{definition}[Vacuum Permittivity Magnitude]
  \label{def:physics:phyx-mini-0983:phyxminiproblems-problemphyxmini0983-vacuumpermittivitymagnitude}
  \lean{PhyXMiniProblems.ProblemPhyXMini0983.VacuumPermittivityMagnitude}
  \uses{def:physics:phyx-mini-0983:phyxminiproblems-problemphyxmini0983-vacuumpermittivitydimension}
  A nonnegative physical vacuum permittivity
\end{definition}
\begin{proof}
  This is the declared model definition of Vacuum Permittivity Magnitude.
\end{proof}
\begin{definition}[Vacuum Permeability Magnitude]
  \label{def:physics:phyx-mini-0983:phyxminiproblems-problemphyxmini0983-vacuumpermeabilitymagnitude}
  \lean{PhyXMiniProblems.ProblemPhyXMini0983.VacuumPermeabilityMagnitude}
  \uses{def:physics:phyx-mini-0983:phyxminiproblems-problemphyxmini0983-vacuumpermeabilitydimension}
  A nonnegative physical vacuum permeability
\end{definition}
\begin{proof}
  This is the declared model definition of Vacuum Permeability Magnitude.
\end{proof}
\begin{definition}[nonnegative SIReadout]
  \label{def:physics:phyx-mini-0983:phyxminiproblems-problemphyxmini0983-nonnegativesireadout}
  \lean{PhyXMiniProblems.ProblemPhyXMini0983.nonnegativeSIReadout}
  Coherent-SI readout of a nonnegative dimensionful quantity
\end{definition}
\begin{proof}
  This is the declared model definition of nonnegative SIReadout.
\end{proof}
\begin{definition}[signed SIReadout]
  \label{def:physics:phyx-mini-0983:phyxminiproblems-problemphyxmini0983-signedsireadout}
  \lean{PhyXMiniProblems.ProblemPhyXMini0983.signedSIReadout}
  Coherent-SI readout of a signed dimensionful quantity
\end{definition}
\begin{proof}
  This is the declared model definition of signed SIReadout.
\end{proof}
\begin{definition}[time In Seconds]
  \label{def:physics:phyx-mini-0983:phyxminiproblems-problemphyxmini0983-timeinseconds}
  \lean{PhyXMiniProblems.ProblemPhyXMini0983.timeInSeconds}
  \uses{def:physics:phyx-mini-0983:phyxminiproblems-problemphyxmini0983-timequantity, def:physics:phyx-mini-0983:phyxminiproblems-problemphyxmini0983-nonnegativesireadout}
  Read physical time in seconds
\end{definition}
\begin{proof}
  This is the declared model definition of time In Seconds.
\end{proof}
\begin{definition}[length In Meters]
  \label{def:physics:phyx-mini-0983:phyxminiproblems-problemphyxmini0983-lengthinmeters}
  \lean{PhyXMiniProblems.ProblemPhyXMini0983.lengthInMeters}
  \uses{def:physics:phyx-mini-0983:phyxminiproblems-problemphyxmini0983-lengthquantity, def:physics:phyx-mini-0983:phyxminiproblems-problemphyxmini0983-nonnegativesireadout}
  Read physical length in metres
\end{definition}
\begin{proof}
  This is the declared model definition of length In Meters.
\end{proof}
\begin{definition}[coefficient In Volts Per Meter Second Squared]
  \label{def:physics:phyx-mini-0983:phyxminiproblems-problemphyxmini0983-coefficientinvoltspermetersecondsquared}
  \lean{PhyXMiniProblems.ProblemPhyXMini0983.coefficientInVoltsPerMeterSecondSquared}
  \uses{def:physics:phyx-mini-0983:phyxminiproblems-problemphyxmini0983-electricfieldquadraticcoefficientmagnitude, def:physics:phyx-mini-0983:phyxminiproblems-problemphyxmini0983-nonnegativesireadout}
  Read eta in volts per metre per second squared
\end{definition}
\begin{proof}
  This is the declared model definition of coefficient In Volts Per Meter Second Squared.
\end{proof}
\begin{definition}[electric Field Component In Volts Per Meter]
  \label{def:physics:phyx-mini-0983:phyxminiproblems-problemphyxmini0983-electricfieldcomponentinvoltspermeter}
  \lean{PhyXMiniProblems.ProblemPhyXMini0983.electricFieldComponentInVoltsPerMeter}
  \uses{def:physics:phyx-mini-0983:phyxminiproblems-problemphyxmini0983-signedelectricfieldcomponent, def:physics:phyx-mini-0983:phyxminiproblems-problemphyxmini0983-signedsireadout}
  Read a signed electric-field component in volts per metre
\end{definition}
\begin{proof}
  This is the declared model definition of electric Field Component In Volts Per Meter.
\end{proof}
\begin{definition}[electric Field Rate In Volts Per Meter Second]
  \label{def:physics:phyx-mini-0983:phyxminiproblems-problemphyxmini0983-electricfieldrateinvoltspermetersecond}
  \lean{PhyXMiniProblems.ProblemPhyXMini0983.electricFieldRateInVoltsPerMeterSecond}
  \uses{def:physics:phyx-mini-0983:phyxminiproblems-problemphyxmini0983-signedelectricfieldrate, def:physics:phyx-mini-0983:phyxminiproblems-problemphyxmini0983-signedsireadout}
  Read a signed electric-field rate in volts per metre per second
\end{definition}
\begin{proof}
  This is the declared model definition of electric Field Rate In Volts Per Meter Second.
\end{proof}
\begin{definition}[magnetic Flux Density Component In Teslas]
  \label{def:physics:phyx-mini-0983:phyxminiproblems-problemphyxmini0983-magneticfluxdensitycomponentinteslas}
  \lean{PhyXMiniProblems.ProblemPhyXMini0983.magneticFluxDensityComponentInTeslas}
  \uses{def:physics:phyx-mini-0983:phyxminiproblems-problemphyxmini0983-signedmagneticfluxdensitycomponent, def:physics:phyx-mini-0983:phyxminiproblems-problemphyxmini0983-signedsireadout}
  Read a signed magnetic-flux-density component in teslas
\end{definition}
\begin{proof}
  This is the declared model definition of magnetic Flux Density Component In Teslas.
\end{proof}
\begin{definition}[magnetic Flux In Webers]
  \label{def:physics:phyx-mini-0983:phyxminiproblems-problemphyxmini0983-magneticfluxinwebers}
  \lean{PhyXMiniProblems.ProblemPhyXMini0983.magneticFluxInWebers}
  \uses{def:physics:phyx-mini-0983:phyxminiproblems-problemphyxmini0983-signedmagneticflux, def:physics:phyx-mini-0983:phyxminiproblems-problemphyxmini0983-signedsireadout}
  Read signed magnetic flux in webers
\end{definition}
\begin{proof}
  This is the declared model definition of magnetic Flux In Webers.
\end{proof}
\begin{definition}[magnetic Flux Rate In Webers Per Second]
  \label{def:physics:phyx-mini-0983:phyxminiproblems-problemphyxmini0983-magneticfluxrateinweberspersecond}
  \lean{PhyXMiniProblems.ProblemPhyXMini0983.magneticFluxRateInWebersPerSecond}
  \uses{def:physics:phyx-mini-0983:phyxminiproblems-problemphyxmini0983-signedmagneticfluxrate, def:physics:phyx-mini-0983:phyxminiproblems-problemphyxmini0983-signedsireadout}
  Read signed magnetic-flux rate in webers per second
\end{definition}
\begin{proof}
  This is the declared model definition of magnetic Flux Rate In Webers Per Second.
\end{proof}
\begin{definition}[electromotive Force In Volts]
  \label{def:physics:phyx-mini-0983:phyxminiproblems-problemphyxmini0983-electromotiveforceinvolts}
  \lean{PhyXMiniProblems.ProblemPhyXMini0983.electromotiveForceInVolts}
  \uses{def:physics:phyx-mini-0983:phyxminiproblems-problemphyxmini0983-electromotiveforce, def:physics:phyx-mini-0983:phyxminiproblems-problemphyxmini0983-signedsireadout}
  Read signed electromotive force in volts
\end{definition}
\begin{proof}
  This is the declared model definition of electromotive Force In Volts.
\end{proof}
\begin{definition}[electric Current In Amperes]
  \label{def:physics:phyx-mini-0983:phyxminiproblems-problemphyxmini0983-electriccurrentinamperes}
  \lean{PhyXMiniProblems.ProblemPhyXMini0983.electricCurrentInAmperes}
  \uses{def:physics:phyx-mini-0983:phyxminiproblems-problemphyxmini0983-electriccurrent, def:physics:phyx-mini-0983:phyxminiproblems-problemphyxmini0983-signedsireadout}
  Read signed electric current in amperes
\end{definition}
\begin{proof}
  This is the declared model definition of electric Current In Amperes.
\end{proof}
\begin{definition}[electric Current Magnitude In Amperes]
  \label{def:physics:phyx-mini-0983:phyxminiproblems-problemphyxmini0983-electriccurrentmagnitudeinamperes}
  \lean{PhyXMiniProblems.ProblemPhyXMini0983.electricCurrentMagnitudeInAmperes}
  \uses{def:physics:phyx-mini-0983:phyxminiproblems-problemphyxmini0983-electriccurrent, def:physics:phyx-mini-0983:phyxminiproblems-problemphyxmini0983-electriccurrentinamperes}
  Read the magnitude of a signed electric current in amperes
\end{definition}
\begin{proof}
  This is the declared model definition of electric Current Magnitude In Amperes.
\end{proof}
\begin{definition}[resistance In Ohms]
  \label{def:physics:phyx-mini-0983:phyxminiproblems-problemphyxmini0983-resistanceinohms}
  \lean{PhyXMiniProblems.ProblemPhyXMini0983.resistanceInOhms}
  \uses{def:physics:phyx-mini-0983:phyxminiproblems-problemphyxmini0983-electricalresistancemagnitude, def:physics:phyx-mini-0983:phyxminiproblems-problemphyxmini0983-nonnegativesireadout}
  Read total circuit resistance in ohms
\end{definition}
\begin{proof}
  This is the declared model definition of resistance In Ohms.
\end{proof}
\begin{definition}[permittivity In Farads Per Meter]
  \label{def:physics:phyx-mini-0983:phyxminiproblems-problemphyxmini0983-permittivityinfaradspermeter}
  \lean{PhyXMiniProblems.ProblemPhyXMini0983.permittivityInFaradsPerMeter}
  \uses{def:physics:phyx-mini-0983:phyxminiproblems-problemphyxmini0983-vacuumpermittivitymagnitude, def:physics:phyx-mini-0983:phyxminiproblems-problemphyxmini0983-nonnegativesireadout}
  Read vacuum permittivity in farads per metre
\end{definition}
\begin{proof}
  This is the declared model definition of permittivity In Farads Per Meter.
\end{proof}
\begin{definition}[permeability In Henries Per Meter]
  \label{def:physics:phyx-mini-0983:phyxminiproblems-problemphyxmini0983-permeabilityinhenriespermeter}
  \lean{PhyXMiniProblems.ProblemPhyXMini0983.permeabilityInHenriesPerMeter}
  \uses{def:physics:phyx-mini-0983:phyxminiproblems-problemphyxmini0983-vacuumpermeabilitymagnitude, def:physics:phyx-mini-0983:phyxminiproblems-problemphyxmini0983-nonnegativesireadout}
  Read vacuum permeability in henries per metre
\end{definition}
\begin{proof}
  This is the declared model definition of permeability In Henries Per Meter.
\end{proof}
\begin{definition}[Figure Feature]
  \label{def:physics:phyx-mini-0983:phyxminiproblems-problemphyxmini0983-figurefeature}
  \lean{PhyXMiniProblems.ProblemPhyXMini0983.FigureFeature}
  Named physical features visible in image 983.png
\end{definition}
\begin{proof}
  This is the declared model definition of Figure Feature.
\end{proof}
\begin{definition}[Figure Label]
  \label{def:physics:phyx-mini-0983:phyxminiproblems-problemphyxmini0983-figurelabel}
  \lean{PhyXMiniProblems.ProblemPhyXMini0983.FigureLabel}
  Symbols printed in the primary image
\end{definition}
\begin{proof}
  This is the declared model definition of Figure Label.
\end{proof}
\begin{definition}[Figure Label Role]
  \label{def:physics:phyx-mini-0983:phyxminiproblems-problemphyxmini0983-figurelabelrole}
  \lean{PhyXMiniProblems.ProblemPhyXMini0983.FigureLabelRole}
  Physical role of each printed label
\end{definition}
\begin{proof}
  This is the declared model definition of Figure Label Role.
\end{proof}
\begin{definition}[Axial Direction]
  \label{def:physics:phyx-mini-0983:phyxminiproblems-problemphyxmini0983-axialdirection}
  \lean{PhyXMiniProblems.ProblemPhyXMini0983.AxialDirection}
  Direction along the vertical cylinder axis
\end{definition}
\begin{proof}
  This is the declared model definition of Axial Direction.
\end{proof}
\begin{definition}[Azimuthal Normal Orientation]
  \label{def:physics:phyx-mini-0983:phyxminiproblems-problemphyxmini0983-azimuthalnormalorientation}
  \lean{PhyXMiniProblems.ProblemPhyXMini0983.AzimuthalNormalOrientation}
  Orientation of the radial--axial rectangle's unit normal
\end{definition}
\begin{proof}
  This is the declared model definition of Azimuthal Normal Orientation.
\end{proof}
\begin{definition}[Current Direction From Above]
  \label{def:physics:phyx-mini-0983:phyxminiproblems-problemphyxmini0983-currentdirectionfromabove}
  \lean{PhyXMiniProblems.ProblemPhyXMini0983.CurrentDirectionFromAbove}
  Current direction resolvable in a view from above the cylinder
\end{definition}
\begin{proof}
  This is the declared model definition of Current Direction From Above.
\end{proof}
\begin{definition}[Diagram Viewpoint]
  \label{def:physics:phyx-mini-0983:phyxminiproblems-problemphyxmini0983-diagramviewpoint}
  \lean{PhyXMiniProblems.ProblemPhyXMini0983.DiagramViewpoint}
  View used by the qualitative direction question
\end{definition}
\begin{proof}
  This is the declared model definition of Diagram Viewpoint.
\end{proof}
\begin{definition}[Maxwell Induction Figure]
  \label{def:physics:phyx-mini-0983:phyxminiproblems-problemphyxmini0983-maxwellinductionfigure}
  \lean{PhyXMiniProblems.ProblemPhyXMini0983.MaxwellInductionFigure}
  \uses{def:physics:phyx-mini-0983:phyxminiproblems-problemphyxmini0983-figurefeature, def:physics:phyx-mini-0983:phyxminiproblems-problemphyxmini0983-figurelabel, def:physics:phyx-mini-0983:phyxminiproblems-problemphyxmini0983-figurelabelrole, def:physics:phyx-mini-0983:phyxminiproblems-problemphyxmini0983-axialdirection, def:physics:phyx-mini-0983:phyxminiproblems-problemphyxmini0983-diagramviewpoint}
  Literal and qualitative evidence transcribed from the primary raster
\end{definition}
\begin{proof}
  This is the declared model definition of Maxwell Induction Figure.
\end{proof}
\begin{definition}[Radial Axial Circuit Geometry]
  \label{def:physics:phyx-mini-0983:phyxminiproblems-problemphyxmini0983-radialaxialcircuitgeometry}
  \lean{PhyXMiniProblems.ProblemPhyXMini0983.RadialAxialCircuitGeometry}
  \uses{def:physics:phyx-mini-0983:phyxminiproblems-problemphyxmini0983-lengthquantity, def:physics:phyx-mini-0983:phyxminiproblems-problemphyxmini0983-azimuthalnormalorientation}
  The conducting circuit and its spanning radial--axial rectangle. The first coordinate of pointAtMeters is radial distance from the cylinder axis and the second is axial distance along the side of length b
\end{definition}
\begin{proof}
  This is the declared model definition of Radial Axial Circuit Geometry.
\end{proof}
\begin{definition}[Displacement Current Induction Setup]
  \label{def:physics:phyx-mini-0983:phyxminiproblems-problemphyxmini0983-displacementcurrentinductionsetup}
  \lean{PhyXMiniProblems.ProblemPhyXMini0983.DisplacementCurrentInductionSetup}
  \uses{def:physics:phyx-mini-0983:phyxminiproblems-problemphyxmini0983-timequantity, def:physics:phyx-mini-0983:phyxminiproblems-problemphyxmini0983-electricfieldquadraticcoefficientmagnitude, def:physics:phyx-mini-0983:phyxminiproblems-problemphyxmini0983-signedelectricfieldcomponent, def:physics:phyx-mini-0983:phyxminiproblems-problemphyxmini0983-signedelectricfieldrate, def:physics:phyx-mini-0983:phyxminiproblems-problemphyxmini0983-signedmagneticfluxdensitycomponent, def:physics:phyx-mini-0983:phyxminiproblems-problemphyxmini0983-signedmagneticflux, def:physics:phyx-mini-0983:phyxminiproblems-problemphyxmini0983-signedmagneticfluxrate, def:physics:phyx-mini-0983:phyxminiproblems-problemphyxmini0983-electromotiveforce, def:physics:phyx-mini-0983:phyxminiproblems-problemphyxmini0983-electriccurrent, def:physics:phyx-mini-0983:phyxminiproblems-problemphyxmini0983-electricalresistancemagnitude, def:physics:phyx-mini-0983:phyxminiproblems-problemphyxmini0983-vacuumpermittivitymagnitude, def:physics:phyx-mini-0983:phyxminiproblems-problemphyxmini0983-vacuumpermeabilitymagnitude, def:physics:phyx-mini-0983:phyxminiproblems-problemphyxmini0983-currentdirectionfromabove, def:physics:phyx-mini-0983:phyxminiproblems-problemphyxmini0983-maxwellinductionfigure, def:physics:phyx-mini-0983:phyxminiproblems-problemphyxmini0983-radialaxialcircuitgeometry}
  Independent physical objects, profiles, and observables. The field profiles are dimensionful scalar components used to calibrate Physlib's spacetime vector fields. The induced current and its direction are independent fields, not definitions made from an answer choice
\end{definition}
\begin{proof}
  This is the declared model definition of Displacement Current Induction Setup.
\end{proof}
\begin{definition}[Matches Problem And Primary Figure]
  \label{def:physics:phyx-mini-0983:phyxminiproblems-problemphyxmini0983-matchesproblemandprimaryfigure}
  \lean{PhyXMiniProblems.ProblemPhyXMini0983.MatchesProblemAndPrimaryFigure}
  \uses{def:physics:phyx-mini-0983:phyxminiproblems-problemphyxmini0983-timeinseconds, def:physics:phyx-mini-0983:phyxminiproblems-problemphyxmini0983-figurefeature, def:physics:phyx-mini-0983:phyxminiproblems-problemphyxmini0983-figurelabel, def:physics:phyx-mini-0983:phyxminiproblems-problemphyxmini0983-displacementcurrentinductionsetup}
  Prose and primary-raster data, containing no induced-current answer
\end{definition}
\begin{proof}
  This is the declared model definition of Matches Problem And Primary Figure.
\end{proof}
\begin{definition}[Satisfies Radial Axial Circuit Geometry]
  \label{def:physics:phyx-mini-0983:phyxminiproblems-problemphyxmini0983-satisfiesradialaxialcircuitgeometry}
  \lean{PhyXMiniProblems.ProblemPhyXMini0983.SatisfiesRadialAxialCircuitGeometry}
  \uses{def:physics:phyx-mini-0983:phyxminiproblems-problemphyxmini0983-lengthinmeters, def:physics:phyx-mini-0983:phyxminiproblems-problemphyxmini0983-displacementcurrentinductionsetup}
  Geometric meaning of the rectangle drawn tight against the axis
\end{definition}
\begin{proof}
  This is the declared model definition of Satisfies Radial Axial Circuit Geometry.
\end{proof}
\begin{definition}[Has Physical Induction Parameters]
  \label{def:physics:phyx-mini-0983:phyxminiproblems-problemphyxmini0983-hasphysicalinductionparameters}
  \lean{PhyXMiniProblems.ProblemPhyXMini0983.HasPhysicalInductionParameters}
  \uses{def:physics:phyx-mini-0983:phyxminiproblems-problemphyxmini0983-timeinseconds, def:physics:phyx-mini-0983:phyxminiproblems-problemphyxmini0983-lengthinmeters, def:physics:phyx-mini-0983:phyxminiproblems-problemphyxmini0983-coefficientinvoltspermetersecondsquared, def:physics:phyx-mini-0983:phyxminiproblems-problemphyxmini0983-resistanceinohms, def:physics:phyx-mini-0983:phyxminiproblems-problemphyxmini0983-permittivityinfaradspermeter, def:physics:phyx-mini-0983:phyxminiproblems-problemphyxmini0983-permeabilityinhenriespermeter, def:physics:phyx-mini-0983:phyxminiproblems-problemphyxmini0983-displacementcurrentinductionsetup}
  Positivity, nondegeneracy, and Physlib calibration assumptions
\end{definition}
\begin{proof}
  This is the declared model definition of Has Physical Induction Parameters.
\end{proof}
\begin{definition}[Satisfies Quadratic Axial Electric Field Model]
  \label{def:physics:phyx-mini-0983:phyxminiproblems-problemphyxmini0983-satisfiesquadraticaxialelectricfieldmodel}
  \lean{PhyXMiniProblems.ProblemPhyXMini0983.SatisfiesQuadraticAxialElectricFieldModel}
  \uses{def:physics:phyx-mini-0983:phyxminiproblems-problemphyxmini0983-coefficientinvoltspermetersecondsquared, def:physics:phyx-mini-0983:phyxminiproblems-problemphyxmini0983-electricfieldcomponentinvoltspermeter, def:physics:phyx-mini-0983:phyxminiproblems-problemphyxmini0983-electricfieldrateinvoltspermetersecond, def:physics:phyx-mini-0983:phyxminiproblems-problemphyxmini0983-displacementcurrentinductionsetup}
  The uniform axial field is calibrated to the Physlib electric field and has the stated quadratic time dependence. Its rate is introduced by a genuine derivative relation rather than by differentiating inside a target formula
\end{definition}
\begin{proof}
  This is the declared model definition of Satisfies Quadratic Axial Electric Field Model.
\end{proof}
\begin{definition}[Satisfies Cylindrical Ampere Maxwell Law]
  \label{def:physics:phyx-mini-0983:phyxminiproblems-problemphyxmini0983-satisfiescylindricalamperemaxwelllaw}
  \lean{PhyXMiniProblems.ProblemPhyXMini0983.SatisfiesCylindricalAmpereMaxwellLaw}
  \uses{def:physics:phyx-mini-0983:phyxminiproblems-problemphyxmini0983-lengthinmeters, def:physics:phyx-mini-0983:phyxminiproblems-problemphyxmini0983-electricfieldrateinvoltspermetersecond, def:physics:phyx-mini-0983:phyxminiproblems-problemphyxmini0983-magneticfluxdensitycomponentinteslas, def:physics:phyx-mini-0983:phyxminiproblems-problemphyxmini0983-permittivityinfaradspermeter, def:physics:phyx-mini-0983:phyxminiproblems-problemphyxmini0983-permeabilityinhenriespermeter, def:physics:phyx-mini-0983:phyxminiproblems-problemphyxmini0983-displacementcurrentinductionsetup}
  Cylindrical Ampere--Maxwell law for circular paths centered on the axis. The right side is vacuum displacement current through the disk of radius r; no requested current in the conducting rectangle appears here
\end{definition}
\begin{proof}
  This is the declared model definition of Satisfies Cylindrical Ampere Maxwell Law.
\end{proof}
\begin{definition}[Satisfies Flux Faraday Ohm And Direction Laws]
  \label{def:physics:phyx-mini-0983:phyxminiproblems-problemphyxmini0983-satisfiesfluxfaradayohmanddirectionlaws}
  \lean{PhyXMiniProblems.ProblemPhyXMini0983.SatisfiesFluxFaradayOhmAndDirectionLaws}
  \uses{def:physics:phyx-mini-0983:phyxminiproblems-problemphyxmini0983-timeinseconds, def:physics:phyx-mini-0983:phyxminiproblems-problemphyxmini0983-lengthinmeters, def:physics:phyx-mini-0983:phyxminiproblems-problemphyxmini0983-magneticfluxdensitycomponentinteslas, def:physics:phyx-mini-0983:phyxminiproblems-problemphyxmini0983-magneticfluxinwebers, def:physics:phyx-mini-0983:phyxminiproblems-problemphyxmini0983-magneticfluxrateinweberspersecond, def:physics:phyx-mini-0983:phyxminiproblems-problemphyxmini0983-electromotiveforceinvolts, def:physics:phyx-mini-0983:phyxminiproblems-problemphyxmini0983-electriccurrentinamperes, def:physics:phyx-mini-0983:phyxminiproblems-problemphyxmini0983-resistanceinohms, def:physics:phyx-mini-0983:phyxminiproblems-problemphyxmini0983-displacementcurrentinductionsetup}
  The aligned field flux is the axial length b times the radial integral from the axis to width a. Faraday--Lenz and passive-sign Ohm law determine the signed emf and current. The direction clauses are generic implications from current sign; they do not assume which sign occurs in this problem
\end{definition}
\begin{proof}
  This is the declared model definition of Satisfies Flux Faraday Ohm And Direction Laws.
\end{proof}
\begin{definition}[Answer Choice]
  \label{def:physics:phyx-mini-0983:phyxminiproblems-problemphyxmini0983-answerchoice}
  \lean{PhyXMiniProblems.ProblemPhyXMini0983.AnswerChoice}
  Labels of the four source expressions
\end{definition}
\begin{proof}
  This is the declared model definition of Answer Choice.
\end{proof}
\begin{definition}[displayed Source Expression]
  \label{def:physics:phyx-mini-0983:phyxminiproblems-problemphyxmini0983-displayedsourceexpression}
  \lean{PhyXMiniProblems.ProblemPhyXMini0983.displayedSourceExpression}
  \uses{def:physics:phyx-mini-0983:phyxminiproblems-problemphyxmini0983-lengthinmeters, def:physics:phyx-mini-0983:phyxminiproblems-problemphyxmini0983-coefficientinvoltspermetersecondsquared, def:physics:phyx-mini-0983:phyxminiproblems-problemphyxmini0983-resistanceinohms, def:physics:phyx-mini-0983:phyxminiproblems-problemphyxmini0983-permittivityinfaradspermeter, def:physics:phyx-mini-0983:phyxminiproblems-problemphyxmini0983-permeabilityinhenriespermeter, def:physics:phyx-mini-0983:phyxminiproblems-problemphyxmini0983-displacementcurrentinductionsetup, def:physics:phyx-mini-0983:phyxminiproblems-problemphyxmini0983-answerchoice}
  Literal scalar transcription of the source choices. These are deliberately not typed as electric currents: as printed, several are dimensionally inconsistent, including choice B's extra factor of the axial length b
\end{definition}
\begin{proof}
  This is the declared model definition of displayed Source Expression.
\end{proof}
\begin{definition}[recorded Dataset Answer]
  \label{def:physics:phyx-mini-0983:phyxminiproblems-problemphyxmini0983-recordeddatasetanswer}
  \lean{PhyXMiniProblems.ProblemPhyXMini0983.recordedDatasetAnswer}
  \uses{def:physics:phyx-mini-0983:phyxminiproblems-problemphyxmini0983-answerchoice}
  The dataset metadata records source choice B
\end{definition}
\begin{proof}
  This is the declared model definition of recorded Dataset Answer.
\end{proof}
% --- Archon declaration topology end ---
% --- Archon physics formalization source end ---
